\section{Homomorfismo}

Consideremos inicialmente a categoria dos anéis associativos comutativos unitários.
Seus objetos são os tais anéis. Entre quaisquer dois tais anéis $A$ e $B$,
consideramos a coleção $\mathrm{Hom}(A,B)$ consistindo em todas as funções
$f : A \to B$, chamadas \underline{homomorfismos} de $A$ em $B$, tais que
\begin{center}
$f(x+y) = f(x) + f(y)$, \\
$f(xy) = f(x)f(y)$, e\\
$f(1) = 1$.
\end{center} 

Para quaisquer tais anéis $A$, $B$ e $C$, notemos duas propriedades fundamentais:

(C1) a composição de funções define uma operação binária associativa
\[
\mathrm{Hom}(A,B) \times \mathrm{Hom}(B,C) \to \mathrm{Hom}(A,C), (f,g) \mapsto gf;
\]

(C2) a aplicação identidade $\mathrm{id}_B \colon B \to B$, $x \mapsto x$,
é um homomorfismo que satisfaz $\mathrm{id}_B \circ f = f$ e $g \circ \mathrm{id}_B = g$
para quaisquer $f \in \mathrm{Hom}(A,B)$ e $g \in \mathrm{Hom}(B,C)$.

\begin{definicao}
Dizemos que dois tais anéis $A$ e $B$ são \underline{isomorfos} se existem
$f \in \mathrm{Hom}(A,B)$ e $g \in \mathrm{Hom}(B,A)$ tais que $gf = \mathrm{id}_A$
e $fg = \mathrm{id}_B$.
\end{definicao}

Por simplicidade, passaremos à categoria dos grupos para demonstrarmos
uma versão multiplicativa mais vigorosa da Proposição 13.2 sobre isomorfismos.
Por definição, um homomorfismo $f : X \to Y$ entre grupos multiplicativos $X$ e $Y$ deve satisfazer
$f(xx') = f(x)f(x')$ para quaisquer $x, x' \in X$.

Fixemos primeiramente as notações de imagem direta e imagem inversa.
Se $f : G \to Z$ é uma função qualquer e $G' \subseteq G$ e
$Z' \subseteq Z$ são subconjuntos de $G$ e $Z$, respectivamente, escrevemos
\[
f(G') = \{f(x) : x \in G'\}
\quad \text{e} \quad
f^{-1}(Z') = \{x \in G : f(x) \in Z'\}.
\]

Quando $f$ é um homomorfismo entre grupos $G$ e $Z$, dizemos que $f(G)$
é a imagem de $f$ e que $f^{-1}(\{e\})$ é o \underline{núcleo} de $f$.
Se $G''$ é também um subconjunto de $G$, denotamos o subconjunto
$\{g'g'' : g' \in G', g'' \in G'' \}$ por $G'G''$.

\begin{proposicao}
Seja $f : G \to Z$ um homomorfismo entre grupos. Sejam $K$ o núcleo de $f$,
$G'$ um subgrupo de $G$ e $Z'$ um subgrupo de $Z$. Valem:
\end{proposicao}

\begin{enumerate}[label=(\alph*), left=0.5cm, align=left, nosep]
    \item Os subconjuntos $f(G')$ e $f^{-1}(Z')$ são subgrupos de $Z$ e $G$,
    respectivamente.

    \item Se $Z'$ está contido em $f(G)$, então a restrição
    $G \to f(G)$ de $f$ induz uma bijeção entre os conjuntos quocientes
    $G/f^{-1}(Z')$ e $f(G)/Z'$.

    \item O subgrupo $K$ é normal em $G$ e os subgrupos
    $f^{-1}(f(G'))$, $KG'$ e $G'K$ coincidem.

    \item Se $N$ é um subgrupo normal de $G'$ que contém $K$, então $f$ induz
    um isomorfismo entre os grupos $G'/N$ e $f(G')/f(N)$.
\end{enumerate}

\begin{proof}
Deixamos as primeiras três afirmações de (a) e (c) como um exercício trivial. Agora, $x \in f^{-1}(f(G'))$ 
se, e somente se, $f(x) = f(y)$ para algum $y \in G'$, e isto equivale a $xy^{-1} = k$ para algum $k \in K$.
Obtemos $f^{-1}(f(G')) = G'K$ de modo semelhante. \\
(b) Como $f(a) \in Z'$ se, e somente se, $a \in f^{-1}(Z')$, a função induzida
$f_* : G/f^{-1}(Z') \to f(G)/Z'$ (como definida na Proposição 13.1) é injetiva. Sendo a dada restrição sobrejetiva,
$f_*$ também o é. \\
(d) Como no argumento do item (b), vemos que a injetividade de $f_* : G'/N \to f(G')/f(N)$ decorre de $f^{-1}(f(N)) = N$, 
pelo item (c). A propriedade de $f_*$ respeitar as operações multiplicativas é herdada de $f$.    
\end{proof}

Por exemplo, se $N$ é um subgrupo normal de um grupo $G$ e $H$ é um subgrupo qualquer de $G$, então a composição 
$H \to HN \to HN/N$ da aplicação inclusão com a projeção canônica induz um isomorfismo $H/(H \cap N) \to HN/N$. 
De fato, basta aplicarmos o Teorema do Isomorfismo (ou a sua generalização no item (d) da Proposição 14.2), já que 
$N$ é um subgrupo normal do grupo $HN$ e a composição $h \mapsto (h \bmod N)$, é um homomorfismo sobrejetivo com núcleo $H \cap N$.

Finalmente, voltemos aos anéis. A imagem inversa de cada ideal por um homomorfismo qualquer também é um ideal. Agora, 
por exemplo, seja $f$ o único elemento em $\mathrm{Hom}(\mathbb{Z}, B)$. (1) Se $B = \mathbb{Q}$, sendo este um corpo, 
ele possui exatamente dois ideais. Então a imagem do ideal $m\mathbb{Z}$ de $\mathbb{Z}$ em $\mathbb{Q}$ só é um ideal 
de $\mathbb{Q}$ se $m = 0$; entretanto, é evidente que $f(m\mathbb{Z})$ é um ideal da imagem $f(\mathbb{Z})$ para qualquer $m$.
(2) Se $B = \mathbb{Z}[i]$ e $J = (2+i)\mathbb{Z}[i]$ é o ideal de $B$ gerado por $2+i$, então $f^{-1}(J)$ é o ideal 
$5\mathbb{Z}$ de $\mathbb{Z}$ e o ideal de $B$ gerado por $f(f^{-1}(J))$ é $5\mathbb{Z}[i]$, o qual está contido propriamente em $J$.

\begin{exemplo}[Subgrupos Normais e Homomorfismos]
\leavevmode
    \begin{enumerate}[left=0.5cm, align=left, nosep]
        \item Dê exemplos de subgrupos normais e relacione-os com homomorfismos.
        
        Recordemos que um subgrupo $H$ de um grupo $G$ é normal se ele satisfaz uma das seguintes condições equivalentes:
        \begin{enumerate}[label=(N\arabic*), left=0.5cm, align=left, nosep]
            \item $ghg^{-1} \in H$ para quaisquer $g \in G$ e $h \in H$;
            \item $ghg^{-1} = H$ para qualquer $g \in G$;
            \item $H$ é o núcleo de um homomorfismo
        \end{enumerate}

        Verifiquemos as equivalências entre essas condições. Que (N2) implica (N1) é imediato, reciprocamente,
        $H = g(g^{-1}Hg)g^{-1} \subseteq gHg^{-1} \subseteq H$. Agora, se $H$ é o núcleo de um homomorfismo, então ele 
        satisfaz (N1), pois $f(ghg^{-1}) = f(g)$ e $f(g)^{-1} = e$. Reciprocamente, se $N$ é um subgrupo normal de $G$, 
        então o núcleo da aplicação canônica $G \to G/H$. Evidentemente, os subgrupos $\{e\}$ e $G$ são normais em $G$.  
        Se $G$ for comutativo, cada subgrupo de $G$ é normal.

        Agora, seja $G$ o grupo $\mathrm{Perm}(X)$, em que $X = \{1,2,3\}$.
        Seja $H$ o subgrupo de $G$ gerado pela permutação
        $\sigma =
        \begin{pmatrix}
        1 & 2 & 3 \\
        2 & 3 & 1
        \end{pmatrix},
        $ de ordem $3$. Verifiquemos que $H$ é um subgrupo normal de $G$ de três modos diferentes.

        (1) O homomorfismo $C_g : G \to G, x \mapsto gxg^{-1}$ é bijetivo, pois $C_g \cdot C_{g^{-1}} = \mathrm{id} = C_{g^{-1}} 
        \cdot C_g$. Como $H$ consiste nos elementos de $G$ que possuem ordem que divide $3$, $C_g(H) \subseteq H$ para cada 
        $g \in G$. \\
        (2) Notemos que as classes laterais à esquerda de $H$ em $G$ são $H$ e $\tau H$ em que 
        $\tau =
        \begin{pmatrix}
        1 & 2 & 3 \\
        2 & 1 & 3
        \end{pmatrix}$. Enquanto isso, as classes laterais à direita são $H$ e $H\tau$.
        Então, $\tau H = H \tau$, ou seja, $\tau H \tau^{-1} = H$. \\
        Sempre que $H$ é um subgrupo com exatamente duas classes laterais, ele é um subgrupo normal.

        (3) Seja $K$ o núcleo do homomorfismo $G \to \mathrm{Perm}(G/H)$, $g \mapsto f(g)$, em que a permutação $f(g): G/H \to G/H$ 
        é dada por $xh \mapsto gxh$. Em geral, podemos descrever $K$ como a interseção dos subgrupos $xHx^{-1}$, com $x \in G$.
        De $K \subseteq H \neq G$ e, pela proposição 13.2, temos as desigualdades $1 < |G/H| \leq |G/K| \leq 2$, existem apenas 
        3 subgrupos normais, logo $K = H$.
        Por fim, como $\mathrm{Perm}(\{1,2,3\})$ é pequeno, não é difícil verificar diretamente que ele possui 
        exatamente $6$ subgrupos.

        Por outro lado, $\mathrm{Perm}(\{1,2,3\})$ é isomorfo a $\mathrm{SL}_2(\mathbb{Z}/2\mathbb{Z})$, (isto é, existe um
        isomorfismo $\sigma \mapsto \begin{pmatrix} 1 & 1 \\ 1 & 0 \end{pmatrix}$ e $\tau \mapsto \begin{pmatrix} 0 & 1 \\ 
        1 & 0 \end{pmatrix}$). Ocorre que, para cada primo $p$, o grupo $G = \mathrm{SL}_2(\mathbb{Z}/p\mathbb{Z})$
        possui o elemento $\begin{pmatrix} -1 & 0 \\ 0 & -1 \end{pmatrix}$, o qual contém todos os elementos de $G$. Daí, o 
        subgrupo $N$ gerado por ele é normal em $G$. Um fato bem conhecido é que, se $p > 2$, então o grupo $G/N$ é simples,
        no sentido de que possui exatamente dois subgrupos normais.

        \item A segunda asserção implica a primeira. Agora, se $z \in \mathbb{C}$, então $z = a + bi$, com $a,b \in \mathbb{R}$ 
        e $i^2 = -1$. Notemos por outro lado, se $f \in A$, então, pela divisão euclidiana, temos que  $f = (X^2 + 1)q + r$,
        com $r = a + bX, \; a,b \in \mathbb{R}$. Donde $(f \bmod I) = (a \bmod I) + (b \bmod I)(X \bmod I)$, com $(X \bmod I)^2 = -1$.
        Consideremos $\psi : \mathbb{C} \to A/I$, $\psi(a+bi) = (a \bmod I) + (b \bmod I)(X \bmod I)$.Vemos que $\psi$ é um 
        homomorfismo de anéis associativos, comutativos e unitários, com núcleo nulo, isto é $\psi$ é um isomorfismo.
    
    \end{enumerate}
\end{exemplo}
